\section{Validation}
\label{sec-validation}

\subsection{Démarche de la campagne}

La campagne de vérification et validation a commencé le 3 octobre 2026 (\cle{docs/VV_campagne.md}). Elle reprend la démarche en deux temps des codes FDEM publiés (Guo et Solidity, Lisjak et Y-Geo~\cite{lisjak2013thesis,mahabadi2012ygeo}, Fukuda, HOSS) : des cas à solution exacte d'abord, la reproduction d'articles ensuite. Elle s'en distingue par une règle : chaque critère d'acceptation est chiffré et écrit avant le premier calcul, et un échec est consigné tel quel, sans retouche du critère. La partie II du document porte les cas à solution exacte (bancs P1 à P5), la partie III la reproduction chiffrée d'articles (B1 à B9), la partie IV un calage unique sur le Red Bohus (compression simple, brésilien, triaxial à 20 et 50~MPa) suivi d'une prédiction sans retouche du triaxial à 75 et 100~MPa.

Chaque banc vit dans un dossier de \cle{vv/} avec un script qui prépare, lance et analyse les calculs, ses critères en constantes, un fichier de résultats et ses figures. Un lanceur répartit les calculs sur les cœurs et reprend les calculs interrompus. Les résultats de la campagne ont été produits sous macOS ARM.

Au 6 octobre 2026, deux bancs ont tourné en campagne (\cref{tab-bancs}). Sept autres sont préparés avec leurs critères et n'ont tourné qu'en essai de fumée, sur maillage grossier ou durée réduite ; leurs chiffres valident la chaîne de calcul, pas le résultat. Les parties III et IV sont vides.

\begin{table}[htbp]
\centering
\caption{Statut des bancs de la campagne de vérification et validation au 6 octobre 2026.}
\label{tab-bancs}
\begin{tabularx}{\linewidth}{@{}l X l X@{}}
\toprule
banc & référence & statut & verdict \\
\midrule
P1.1 onde dans une barre & d'Alembert & exécuté & \cle{fem3d} passe ; FDEM adaptatif échoue au seul bilan ; intrinsèque $p_f = 20$ échoue \\
P2.1 bloc glissant & $L = v_0^2/(2\mu g)$ & exécuté & potentiel passe ; pénalité échoue \\
P1.2, P1.3 élasticité & Kirsch, champ homogène & fumée & aucun \\
P2.2 sphère sur un plan & Hertz & fumée & aucun \\
P3.1 fissure pressurisée & mécanique de la rupture, Guo & fumée & aucun \\
P3.3 fissure en aile & Lisjak, Y-Geo & fumée & aucun \\
P4.1, P4.2 flexion, brésilien & poutre, Hondros, Guo & essai à blanc & aucun \\
B4 à B7 articles & AbuAisha, Wang, Yan, Yang & B6 relancé & B6 : reproduction en échec \\
B8, B9 impact et coupe & Saksala, Heilman, LeBaron & fumée & aucun \\
P3.2 énergie de fissure & $G_f$ fois l'aire rompue & non préparé & aucun \\
\bottomrule
\end{tabularx}
\end{table}

\subsection{Onde de compression dans une barre}
\label{sec-V1}

Ce banc vérifie la célérité, l'impédance et les réflexions, et mesure un ordre de convergence : une erreur à ce niveau se retrouverait dans tous les bancs d'impact. La barre mesure $8 \times 8 \times 400$~mm, avec $\rho = 2\,650$~kg/m$^3$, $E = 60$~GPa et $\nu = 0{,}25$ ; ses quatre faces latérales sont sur rouleaux, ce qui rend le problème exactement unidimensionnel, sans la dispersion de Pochhammer, avec $c = \sqrt{M/\rho} = 5\,212{,}47$~m/s. Une vitesse de $0{,}1$~m/s en trapèze de 4, 12 et $4~\mu$s est imposée à une extrémité, l'autre est libre, et le calcul dure $250~\mu$s, soit $1{,}6$ aller-retour. La solution exacte est la série de d'Alembert avec réflexions ; le plateau de la réaction vaut $\rho c v_0 A = 88{,}40$~N. Les critères, fixés avant calcul sur le maillage le plus fin, sont une erreur $L^2$ du déplacement d'au plus $2~\%$, des erreurs de célérité et de force d'au plus 1 et $2~\%$, un résidu du bilan d'au plus $10^{-6}~\%$ et un ordre de convergence d'au moins 1.

\begin{table}[htbp]
\centering
\caption{Onde dans une barre : erreurs mesurées par variante et taille d'élément $h$.}
\label{tab-V1}
\small
\begin{tabular}{@{}lrrrrrr@{}}
\toprule
variante & $h$ (mm) & tétraèdres & $e_u$ & $c$ & $F$ & résidu du bilan \\
\midrule
\cle{fem3d} & 4 & 3\,688 & $0{,}46~\%$ & $-0{,}05~\%$ & $-0{,}04~\%$ & $2{,}6\times10^{-12}~\%$ \\
\cle{fem3d} & 1 & 125\,486 & $0{,}05~\%$ & $<0{,}01~\%$ & $<0{,}01~\%$ & $1{,}2\times10^{-10}~\%$ \\
FDEM adaptatif & 2 & 20\,435 & $0{,}17~\%$ & $+0{,}04~\%$ & $<0{,}01~\%$ & $5{,}5\times10^{-4}~\%$ \\
FDEM intrinsèque $p_f = 20$ & 4 & 3\,688 & $7{,}7~\%$ & $-2{,}9~\%$ & $-2{,}9~\%$ & $1{,}5\times10^{-3}~\%$ \\
FDEM intrinsèque $p_f = 20$ & 2 & 20\,435 & $8{,}6~\%$ & $-3{,}2~\%$ & $-3{,}3~\%$ & $3{,}5\times10^{-4}~\%$ \\
\bottomrule
\end{tabular}
\end{table}

Le solveur éléments finis passe les cinq critères (\cref{tab-V1}), avec un ordre observé de $1{,}58$ sur le déplacement (\cref{fig-V1conv}). L'ordre inférieur à 2 est attribué aux quatre angles du trapèze, qui limitent la régularité de la solution exacte. L'erreur $L^2$ sur la réaction décroît de $16~\%$ à $3{,}5~\%$ ; elle vient des oscillations de dispersion numérique au retour de l'onde, le plateau étant exact.

La FDEM en insertion adaptative reproduit les éléments finis à $10^{-4}$ près maillage par maillage, avec un ordre de $1{,}46$, mais échoue au critère du bilan : son résidu vaut $2{,}2\times10^{-3}$ à $5{,}5\times10^{-4}~\%$. Le document l'attribue au décalage d'un demi-pas entre les compteurs de travail et le saute-mouton, résidu d'ordre $\Delta t$ qui décroît avec le pas, et non à une création d'énergie. Le critère, calibré sur les éléments finis, est jugé inadapté mais n'a pas été changé après coup.

\begin{figure}[htbp]
\centering
\includegraphics[width=0.62\linewidth]{fig_convergence}
\caption{Onde dans une barre : erreur $L^2$ du déplacement à mi-longueur contre la taille d'élément, pour les éléments finis, la FDEM adaptative et la FDEM intrinsèque à $p_f = 20$.}
\label{fig-V1conv}
\end{figure}

La FDEM en insertion intrinsèque à $p_f = 20$, réglage du banc d'impact, échoue : l'onde est $2{,}9$ à $3{,}2~\%$ trop lente et le plateau de force trop faible d'autant, sans aucune rupture, et l'erreur ne décroît pas avec $h$ (ordre $-0{,}16$), parce que le nombre de joints par unité de longueur croît comme $1/h$. Avec ce réglage, les ondes d'un impact arrivent environ $3~\%$ en retard et la force de contact initiale est sous-estimée d'environ $3~\%$, indépendamment du maillage.

Un balayage de la pénalité de 10 à 200 a testé un modèle de joints en série, $c_\text{eff}/c = (1 + \alpha/p_f)^{-1/2}$, avec $\alpha = 1{,}21$ calé sur le seul point $p_f = 20$ et écrit avant le balayage (\cref{fig-V1pen}). La prédiction tient à $0{,}05$ point près sur toute la plage ; l'ajustement sur les huit points donne $\alpha = 1{,}239$. Le retard vaut $-5{,}54~\%$ à $p_f = 10$, $-1{,}23~\%$ à 50, $-0{,}65~\%$ à 100 et $-0{,}35~\%$ à 200, ce dernier indépendant de $h$. Le pas de temps décroît comme $1/\sqrt{p_f}$. La règle qui en sort est $p_f \geq \alpha/(2\varepsilon) \approx 0{,}62/\varepsilon$ pour un biais de célérité $\varepsilon$ : 62 pour $1~\%$, 124 pour $0{,}5~\%$, au prix d'un pas plus petit d'un facteur $1{,}8$ ou $2{,}5$ par rapport à $p_f = 20$. Le banc de la fissure pressurisée en tient déjà compte avec $p_f = 100$.

\begin{figure}[htbp]
\centering
\includegraphics[width=0.62\linewidth]{fig_penalite}
\caption{Retard de célérité de l'onde contre le facteur de pénalité intrinsèque, mesuré sur deux maillages, avec la prédiction écrite avant le balayage et l'ajustement.}
\label{fig-V1pen}
\end{figure}

Le banc a trois réserves. La barre est confinée latéralement, la barre libre dispersive n'est pas couverte. Le pas de la FDEM est commandé par les tétraèdres aplatis (hauteur inscrite de $0{,}22$~mm à $h = 2$~mm), ce qui rend le calcul 40 à 100 fois plus coûteux qu'en éléments finis. Les résultats dépendent du générateur de maillage au-delà de la quatrième décimale.

\subsection{Bloc glissant jusqu'à l'arrêt}
\label{sec-P21}

Ce banc vérifie le frottement de Coulomb du contact général, d'après Xiang et al. 2009, repris en 3D par Fukuda et al. Un cube de 50~mm ($E = 1$~GPa) glisse sur une dalle fixe avec $\mu = 0{,}5$ sous la pesanteur ; les deux corps sont continus, leurs maillages ne coïncident pas, et seul le contact général les relie. Sans basculement, le bloc s'arrête à $L = v_0^2/(2\mu g)$. Les critères sont un écart d'au plus $2~\%$ sur la distance d'arrêt et sur la trajectoire.

\begin{table}[htbp]
\centering
\caption{Bloc glissant : distance d'arrêt simulée contre exacte selon la loi de contact et la vitesse initiale.}
\label{tab-P21}
\small
\begin{tabular}{@{}lrrrrl@{}}
\toprule
contact & $v_0$ (m/s) & $L$ simulé (mm) & $L$ exact (mm) & écart & verdict \\
\midrule
potentiel & $0{,}5$ & $25{,}93$ & $25{,}48$ & $+1{,}74~\%$ & passe \\
potentiel & 1 & $102{,}92$ & $101{,}94$ & $+0{,}96~\%$ & passe \\
potentiel & 2 & $410{,}27$ & $407{,}75$ & $+0{,}62~\%$ & passe \\
pénalité & $0{,}5$ & $-0{,}46$ & $25{,}48$ & $-102~\%$ & échoue \\
\bottomrule
\end{tabular}
\end{table}

Le contact par potentiel reproduit la cinématique aux trois vitesses (\cref{tab-P21,fig-P21}). Le bloc glisse un peu trop loin, d'autant moins que la vitesse est grande, ce que le document attribue à la régularisation du frottement à très faible vitesse en fin d'essai. Sur le premier millimètre, le travail de frottement égale $\mu m g$ fois la distance à $10^{-5}$ près. Le contact par pénalité, défaut de rockim, échoue : le bloc accroche par son arête les triangles de la dalle, bascule, son centre monte de $4{,}7$~mm, il retombe en arrière et finit $0{,}46$~mm derrière son point de départ, avec un bilan d'énergie pourtant fermé à $7\times10^{-12}~\%$. La règle qui en sort impose le contact par potentiel dès qu'un corps glisse sur un autre à travers des maillages non coïncidents ; les configurations d'impact l'appliquent déjà. Le basculement n'est jugé qu'indirectement : le critère de rotation déclaré dans le script n'est pas calculé.

\begin{figure}[htbp]
\centering
\includegraphics[width=0.62\linewidth]{fig_bloc}
\caption{Bloc glissant : déplacement contre temps pour trois vitesses initiales, contact par potentiel et par pénalité, comparé à la solution exacte en trait large.}
\label{fig-P21}
\end{figure}

\subsection{Bancs préparés}

Les bancs suivants ont leurs critères écrits et un essai de fumée ; aucun n'a de résultat de campagne. Le banc d'élasticité P1.2-P1.3 vérifie un champ homogène à la précision machine et la solution de Kirsch autour d'un trou ; l'essai de fumée passe le champ homogène à $7{,}5\times10^{-8}$. Le banc de la sphère sur un plan P2.2 compare le contact par potentiel à Hertz sur quatre raffinements ; l'essai court donne une force filtrée à $+4{,}6~\%$ et un enfoncement à $-3{,}6~\%$ de Hertz, le pic brut étant à $+36~\%$. Le banc de la fissure pressurisée P3.1 reprend le cas de Guo et doit quantifier l'objectivité de la rupture vis-à-vis du maillage, lacune ouverte du code ; la variante la plus fine est estimée à 71 à 114~h de calcul. Le banc de la fissure en aile P3.3 vise la contrainte d'amorçage de $7{,}0$~MPa et l'angle de $64^\circ$ de Lisjak ; deux des quatre formules analytiques de comparaison sont des reconstructions non vérifiées. Les bancs de flexion et de brésilien P4 reprennent Guo ; les critères écrits avant calcul annoncent qu'ils ne peuvent pas tous passer, parce que Guo trouve une résistance brésilienne de $0{,}65\,f_t$ alors qu'un modèle qui respecte Hondros doit amorcer au centre près de $f_t$. Le banc B8 compare le portage de la loi de Saksala à l'impact de bouton publié~\cite{saksala2011ijnamg} ; c'est une comparaison de code à code avec la même loi.

\subsection{Reproduction de Yang et al. 2025 : calcaire de St Anne}
\label{sec-stanne}

L'impact d'un insert sur le calcaire de St Anne à $10{,}66$~m/s~\cite{yang2025validation} est la comparaison la plus aboutie, et la seule menée sur un calcul complet de $300~\mu$s. Le maillage compte 108\,667 tétraèdres, dont un cœur de 14\,030 à $1{,}42$~mm d'arête médiane sous l'insert ; le matériau est celui de la table~4 de l'article ; la pénalité est prise sur les arêtes avec $p_0 = 3\,000$~GPa ; la loi de joint est \cle{plastic} avec plage de Coulomb et cliquet ; la pulvérisation et la viscosité sont désactivées. Le calcul a duré 29~h~21 sur 14 fils.

Les sept critères de l'article ne sont publiés que sous forme de graphiques et ont été lus à $10~\%$ près au mieux ; la chronologie vient du texte et est exacte (\cref{tab-stanne}). La contrainte au taillant, la vitesse d'indentation, le rayon du cratère et la fin de la phase de charge tombent à 6 à $12~\%$ de la simulation de Yang et al., de l'ordre de la dispersion expérimentale et du coefficient de variation de $10~\%$ annoncé par les auteurs. L'enfoncement et la masse de fragments s'écartent nettement, dans le même sens d'un cratère trop ouvert ; trois causes possibles ne sont pas départagées : le réglage du brossage des fragments, un cratère réellement plus ouvert, des éléments comptés qu'un brossage physique ne détacherait pas. La longueur des fissures radiales dépend d'une convention non tranchée, avec un facteur $1{,}7$ entre les deux définitions possibles. Le rebond n'est pas mesuré, car il faudrait au moins $450~\mu$s : la configuration à $500~\mu$s est prête mais n'a pas été lancée.

\begin{table}[htbp]
\centering
\caption{Réplique St Anne à $300~\mu$s : critères de Yang et al. 2025 lus sur les figures 9 et 10, et valeurs de rockim.}
\label{tab-stanne}
\small
\begin{tabular}{@{}lrrrr@{}}
\toprule
critère & Yang FDEM & Yang essai & rockim & écart \\
\midrule
contrainte au taillant & 200~MPa & 200~MPa & $212{,}3$~MPa & $+6~\%$ \\
vitesse d'indentation & $8{,}0$~m/s & $7{,}9$~m/s & $7{,}34$~m/s & $-8~\%$ \\
enfoncement maximal & $1{,}10$~mm & $1{,}05$~mm & $1{,}292$~mm & $+17~\%$ \\
vitesse de rebond & $5{,}6$~m/s & $4{,}2$~m/s & non mesurée & sans objet \\
longueur radiale & 32~mm & 31~mm & $29{,}6$ ou $17{,}7$~mm & $-8$ ou $-45~\%$ \\
rayon du cratère & $13{,}5$~mm & 13~mm & $12{,}00$~mm & $-11~\%$ \\
masse de fragments & $2{,}5$~g & $1{,}2$~g & $4{,}70$~g & $+88~\%$ \\
fin de la phase de charge & $254~\mu$s & non publiée & $284{,}3~\mu$s & $+12~\%$ \\
\bottomrule
\end{tabular}
\end{table}

Le bilan d'énergie du calcul ferme : l'énergie cinétique passe de $60{,}12$ à $3{,}28$~J, le résidu vaut $-9{,}8\times10^{-8}$~J, soit $1{,}5\times10^{-7}~\%$ de l'échelle, sur 12\,090 facettes rompues. La rupture dissipe $34{,}5$~J, le frottement $27{,}8$~J, les joints $14{,}5$~J, et $18{,}1$~J restent stockés en énergie élastique. Le réseau de fissures forme un noyau broyé sous l'insert et des branches radiales (\cref{fig-stanne-joints}) ; la partition entre traction et cisaillement des couleurs n'est qu'indicative, à cause de la normalisation retenue pour le mode de rupture.

\begin{figure}[htbp]
\centering
\includegraphics[width=\linewidth]{stanne_joints}
\caption{Réplique St Anne à $300~\mu$s : 12\,090 facettes rompues en vue 3D, de dessus et en coupe, colorées par mode de rupture (haut) et par type, fissure ou broyage (bas).}
\label{fig-stanne-joints}
\end{figure}

La force de contact mesurée entre l'insert et la roche croît régulièrement jusqu'à environ 90~kN à la fin de la charge (\cref{fig-stanne-fp}). L'estimateur cinématique du train rigide, employé dans les comparaisons antérieures, oscille autour de cette courbe et surestime son maximum de $90~\%$, alors que les deux impulsions cumulées coïncident à $3{,}4~\%$.

\begin{figure}[htbp]
\centering
\includegraphics[width=\linewidth]{stanne_fp}
\caption{Réplique St Anne : force de contact insert-roche mesurée contre l'estimateur du train rigide, en fonction de la pénétration et du temps, impulsions cumulées et autres paires de contact.}
\label{fig-stanne-fp}
\end{figure}

Cette comparaison n'est pas indépendante du chemin de mise au point. La loi de joint et la pénalité ont été choisies pendant la même campagne, en regardant les critères de Yang et al. (\cle{vv/B_articles/README.md}, l.~256-260). Une seule résolution de maillage a tourné ; le maillage plus fin (251\,460 tétraèdres) reste à jouer. Les valeurs de référence elles-mêmes ne sont pas stables d'un document à l'autre : le bilan d'août citait une table de l'article avec une indentation de $9{,}40$ à $9{,}85$~m/s et un enfoncement de $1{,}53$~mm, que la comparaison de septembre, lue sur les figures 9 et 10, ne retrouve pas.

\subsection{Reproduction de Yang et al. 2026 : granite de Kuru}

L'impact sur le granite de Kuru à 9~m/s avec pulvérisation~\cite{yang2026pulverisation} était la cible initiale ; il a été délaissé au profit de St Anne parce que le modèle de pulvérisation n'est pas complètement publié. Il n'est pas reproduit (\cref{tab-kuru-ecarts}). La vitesse d'indentation est trop grande de $22~\%$, la pulvérisation quasi inactive (2 éléments contre environ 360), le freinage moyen trop faible, et le calcul s'arrête avant le retournement du bit.

\begin{table}[htbp]
\centering
\caption{Réplique Kuru, maillage $s = 1$ : grandeurs de rockim et de Yang et al. 2026.}
\label{tab-kuru-ecarts}
\begin{tabular}{@{}lrr@{}}
\toprule
grandeur & rockim & Yang et al. 2026 \\
\midrule
vitesse d'indentation, pente 10-$90~\%$ & $6{,}86$~m/s & $5{,}62$~m/s \\
compression maximale à la jauge & $175{,}95$~MPa & environ 160~MPa \\
éléments pulvérisés & 2 à $182{,}9~\mu$s & environ 360 \\
tétraèdres & 120\,185 & 230\,788 \\
masse du piston & $1{,}057$~kg & $1{,}173$~kg \\
freinage moyen & 20 à 25~kN (estimateur) & environ 57~kN (déduit) \\
\bottomrule
\end{tabular}
\end{table}

Les explications documentées sont partielles. La branche tangentielle de la loi de joint active n'était pas celle de Guo (pente $p_j$ au lieu de $2p_j$). La pulvérisation dépend d'une définition de la mesure de déformation que l'article ne donne pas. Les masses du train sont $10~\%$ plus faibles que chez Yang et al. Le freinage est comparé entre un maximum d'estimateur biaisé et une moyenne déduite. La chaîne causale qui irait d'un lit de débris trop lâche à une réaction faible, puis à l'absence de pulvérisation et de fissures radiales, est donnée par le document comme une hypothèse cohérente avec les données et non comme une causalité démontrée. Un diagnostic indépendant du 12 septembre conclut que les données ne démontrent pas une incapacité de la FDEM 3D à localiser la rupture, mais une reproduction encore non validée ; il a corrigé deux défauts de dépouillement (10~\% de faux positifs dans le décompte des joints rompus, rayon maximal pris pour une longueur de fissure).

La transcription mot à mot de la loi de joint de Solidity a établi que cette loi crée de l'énergie sur cycle fermé (\cref{sec-decharge}). Sur l'impact sans garde-fou, le travail cumulé des joints atteint 2\,745~J pour $31{,}5$~J apportés. Cela n'établit pas que la version interne du code de Yang et al. diverge de la même façon. Le code public de Solidity ne contient plus rien qui n'ait été transcrit : l'effet de vitesse interne, la pulvérisation, la configuration du granite et l'amortissement n'y figurent pas. Une demande de données aux auteurs est rédigée et n'a pas été envoyée.

\subsection{Autres comparaisons}

Sur la fracturation hydraulique en paroi de forage d'AbuAisha et al.~\cite{abuaisha2017hm}, calculée en 2D avec un couplage hydromécanique sans écoulement, la pression de rupture est trop forte de $+13{,}2~\%$ en isotrope et de $+24{,}9$ à $+25{,}0~\%$ en anisotrope ; l'ouverture d'une fissure mûre est à $-4~\%$ de Sneddon et le déplacement de paroi à $-2{,}0~\%$ de Lamé. Sur la compression simple 2D avec insertion adaptative de Yan et al.~\cite{yan2023adaptive}, la résistance de 51~MPa est retrouvée, ce qui reproduit leur simulation calée sans rien valider d'indépendant ; relancée avec le binaire courant, elle dérive de $+0{,}70~\%$ pour une tolérance de $0{,}29~\%$, avec $34~\%$ de joints rompus en plus, sans cause départagée entre l'évolution du code et le maillage de Voronoï. Sur le tunnel de Wang et al., le rayon de la zone endommagée vaut $17{,}19$~m contre 19~m. Ces comparaisons portent sur la simulation calée d'un autre code ou sur un critère analytique ; aucune n'est une validation au sens de la partie IV.

\subsection{Ce qui n'est pas validé}

Aucune comparaison n'a été faite à un essai qui n'a pas servi au calage, aucun essai dynamique de laboratoire (barres de Hopkinson, brésilien dynamique) n'a été confronté, et les données d'essai d'Aising et Gerbaud, base expérimentale des articles de Yang et al. et disponibles à Mines Paris-PSL, ne sont pas exploitées. Le modèle de pulvérisation n'est pas identifié. La loi \cle{plastic} n'est pas éprouvée hors des trajets testés. La convergence en maillage de la fissuration n'a qu'une résolution exploitable sur St Anne, et l'objectivité de la rupture vis-à-vis du maillage n'est pas quantifiée. Le relais du joint au contact n'est pas validé sur l'impact. L'énergie dissipée par la fissuration n'a jamais été comparée à $G_f$ multiplié par l'aire rompue.

\figaproduire{4.5cm}{écart relatif de chaque grandeur de comparaison (St Anne, Kuru, AbuAisha, Yan) à sa référence, avec le statut de la référence (simulation calée, essai, analytique) et la bande de dispersion expérimentale.}{aucun calcul nouveau ; agrégation des fichiers de résultats de \cle{vv/} et des tableaux de ce rapport, figure annoncée par \cle{vv/B_articles/README.md} mais absente.}
